\chapter*{怎么读这本手册}
\addcontentsline{toc}{chapter}{怎么读这本手册}

\begin{beginnerbox}[你需要的背景]
\begin{itemize}[nosep]
  \item 会看 \keyterm{if/for/数组} 这种最基础的代码即可；
  \item 不需要懂备份引擎、不需要懂数学系拓扑；
  \item 每个新概念都会先用\textbf{生活比喻}，再用\textbf{小图}，最后才给\textbf{一行伪代码}。
\end{itemize}
\end{beginnerbox}

\section*{两代算法指什么}

\begin{table}[htbp]
\centering
\caption{本手册的「两代」}
\begin{tabular}{@{}cll@{}}
\toprule
\textbf{代} & \textbf{名字} & \textbf{一句话} \\
\midrule
第一代 & FastCDC Stream & 默认切法：一个滚动指纹 + 一个「路标」 \\
第二代 & G-TCDC v6 2F-Gear & 新切法：两个滚动指纹 + 两个「路标」 \\
\bottomrule
\end{tabular}
\end{table}

\section*{颜色约定}

\begin{itemize}[nosep]
  \item \textcolor{Gen1Blue}{蓝色} = 第一代 FastCDC
  \item \textcolor{Gen2Teal}{青绿色} = 第二代 2F-Gear
  \item \textcolor{CutRed}{红色虚线} = 切点（chunk 结束）
\end{itemize}

\section*{图示约定}

\begin{figure}[htbp]
\centering
\begin{tikzpicture}[node distance=0.5cm]
  \node[byte] (b) {字节格};
  \node[gen1, right=0.6cm of b, minimum width=2cm] (g1) {第一代框};
  \node[gen2, right=0.6cm of g1, minimum width=2cm] (g2) {第二代框};
  \node[zone=CutRed, right=0.6cm of g2, minimum width=2cm] (z) {护栏/警告};
  \draw[cutline] ([yshift=-1.2cm]b.south west) -- ++(0,0.8)
    node[above,font=\scriptsize,xshift=2.5cm] {红色虚线 = 切点};
\end{tikzpicture}
\caption{全书统一视觉语言，看图即可分辨「哪一代、是不是切点」。}
\end{figure}

\section*{建议阅读顺序}

按章节顺序读即可。若只想快速知道「第二代新在哪」，可直接跳到第 \ref{ch:gen2-big} 章，再回读第一代第 \ref{ch:gen1-big} 章对照。

\begin{simbox}[全书统一模拟文件]
各章「模拟示例」与第 9、14 章完整走查使用\textbf{同一份 12 字节玩具文件}和\textbf{同一张 gear/norm 表}（见 \filepath{chapters/sim-toy-data.tex}）。\\
建议：先扫每章模拟框 $\to$ 再读第 9 章逐步表 $\to$ 最后对照第 14 章看两代切点差异。
\end{simbox}

\begin{figure}[htbp]
\centering
\begin{tikzpicture}[node distance=0.3cm]
  \node[gen1, minimum width=2cm] (a) {Part I\\动机};
  \node[gen1, minimum width=2cm, right=0.3cm of a] (b) {Part II\\第一代};
  \node[gen2, minimum width=2cm, right=0.3cm of b] (c) {Part III\\第二代};
  \node[gen2, minimum width=2cm, right=0.3cm of c] (d) {Part IV\\对照};
  \draw[arr] (a) -- (b) -- (c) -- (d);
\end{tikzpicture}
\caption{约 45--55 页；每章 2--4 张图，可一天读完 Part II。}
\end{figure}

\vspace{1em}
\noindent\textit{本手册只讲「怎么切 chunk」，不讲压缩、加密、仓库格式。}
